\documentclass{article}
\usepackage{amsmath, amssymb, amsthm}
\usepackage{hyperref}
\usepackage{geometry}
\geometry{margin=1in}

\title{Erd\H{o}s Problem \#655}
\date{Last updated: 2025-08-31}
\author{Source: erdosproblems.com}

\begin{document}
\maketitle

\noindent\textbf{Status:} OPEN \quad \textbf{No prize}

\noindent\textbf{Tags:} geometry, distances

\medskip
\noindent\textbf{Problem Statement:}

Let $x_1,\ldots,x_n\in \mathbb{R}^2$ be such that no circle whose centre is one of the $x_i$ contains three other points. Are there at least\[(1+c)\frac{n}{2}\]distinct distances determined between the $x_i$, for some constant $c>0$ and all $n$ sufficiently large?

\bigskip
\noindent\textbf{Remarks:}

A problem of Erd\H{o}s and Pach. It is easy to see that this assumption implies that there are at least $\frac{n-1}{2}$ distinct distances determined by every point.

Zach Hunter has observed that taking $n$ points equally spaced on a circle disproves this conjecture. In the spirit of related conjectures of Erd\H{o}s and others, presumably some kind of assumption that the points are in general position (e.g. no three on a line and no four on a circle) was intended.

\bigskip
\noindent\small{Source: \url{https://www.erdosproblems.com/655}}
\end{document}
